\scanpage{40}\mspage{5}
% archive sheet number: 39
On appellera \uline{orientation} de $(K,(D_i)_{i\in J})$, la
donnée, pour $\forall\, I\in\unsure{\mathfrak P_3|\mathcal P_3}(J)$ comme dessus
(de $\forall$ ens. de 3 pseudo-droites non concourantes),
d'une \uline{orientation} \inserted{$\omega_I$} du graphe
octaédral gauche correspondant.

Soit maintenant $\mathcal X$ une \uline{carte} qui
soit un arrangement de pseudo-droites.
Son 1-squelette \inserted{$K$} est alors muni de \deleted{une}\inserted{\unsure{$\mathcal C$|la}}
famille $(D_i)_{i\in J}$ de pseudo-droites, c'est
donc un arrangement $\underline{1}$ dim. de
ps. droites. De plus il est muni d'
une orientation naturelle $\omega=(\omega_I)_{\substack{I\in\mathfrak P_3(J)\\\bigcap_{i\in I}D_i=\varnothing}}$.

On trouve ainsi un foncteur
\[\mathcal X\longmapsto (K,(D_i)_{i\in J},\omega)\tag*{(18)}\]
de \unsure{lá|la} catégorie des arrangem\textsuperscript{ts} comb. de
pseudo-droites, vers la catégorie combinatoire
des arrangem\textsuperscript{ts} $1$ dimensionnels orientés de
pseudo-droites.

% slanted
\margin{\deleted{\unsure{Il faut restreindre} le foncteur aux} Il est vrai que \unsure{si l'on} se limite aux \unsure{composantes|comp.} \unsure{comp. part.} qui ne \unsure{corr. pas} \unsure{de} \unsure{trigones|trigones et} \unsure{bigônes|bigones|bigone}}%
\uline{\unsure{Prop\textsuperscript{n}.|Prop.|Th.}} Le foncteur précédent est \uline{pleinement}
\uline{fidèle}, i.e. pour deux arrangem\textsuperscript{ts} comb.
$\mathcal X,\mathcal X'$ de ps. droites, on a
\[\operatorname{Isom}(\mathcal X,\mathcal X')\xrightarrow{\sim}\operatorname{Isom}\big((K,(D_i)_{i\in J},\omega),(K',(D'_{i'})_{i'\in J'},\omega')\big)\]
